% part: intuitionistic-logic
% chapter: introduction
% section: constructive-reasoning

\documentclass[../../../include/open-logic-section]{subfiles}

\begin{document}

\olfileid{int}{int}{cr}

\olsection{استدلال سازنده}

منطق شهودگرایانه، برخلاف گسترش‌های منطق کلاسیک با عملگرهای موجهاتی یا
سورهای مرتبهٔ دوم، از آن رو «ناکلاسیک» است که منطق کلاسیک را محدود
می‌کند. منطق کلاسیک از جهات گوناگون \emph{ناسازنده} است. هدف منطق
شهودگرایانه آن است که گونه‌ای «سازنده‌تر» از استدلال را صورت‌بندی کند
که مشخصهٔ نوعی ریاضیات سازنده است. مثال‌های زیر شاید برخی انگیزه‌های
بنیادی را روشن کنند.

فرض کنید کسی ادعا کند عدد طبیعی~$n$ای یافته است با این خاصیت که اگر
$n$ زوج باشد، فرضیهٔ ریمان صادق است، و اگر~$n$ فرد باشد، فرضیهٔ ریمان
کاذب است. چه خبر خوبی!{} صدق یا کذب فرضیهٔ ریمان یکی از پرسش‌های بزرگ
گشودهٔ ریاضیات است و ظاهراً آن شخص مسئله را به یک محاسبه فروکاسته است:
تعیین اینکه عددی معین زوج است یا نه.

مقدار جادویی~$n$ چیست؟ او چنین توصیفش می‌کند: $n$ عدد طبیعی‌ای است که
اگر فرضیهٔ ریمان صادق باشد برابر~$2$ است، وگرنه برابر~$3$.

خشمگینانه پول خود را پس می‌خواهید. از دیدگاه کلاسیک، توصیف بالا واقعاً
مقداری یگانه برای~$n$ تعیین می‌کند؛ اما آنچه حقیقتاً می‌خواستید مقداری
از~$n$ بود که \emph{به‌صراحت} داده شود.

برای نمونه‌ای دیگر، که شاید کم‌تر ساختگی باشد، پرسش زیر را در نظر
بگیرید. می‌دانیم می‌توان عددی گنگ را به توانی گویا رساند و نتیجه‌ای گویا
گرفت؛ برای نمونه، $\sqrt{2}^2 = 2$. آنچه کم‌تر روشن است این است که آیا
می‌توان عددی گنگ را به توانی \emph{گنگ} رساند و نتیجه‌ای گویا گرفت.
قضیهٔ زیر پاسخ مثبت می‌دهد:

\begin{thm}
اعداد گنگ~$a$ و~$b$ای وجود دارند که~$a^b$ گویاست.
\end{thm}

\begin{proof}
$\sqrt{2}^{\sqrt{2}}$ را در نظر بگیرید. اگر این عدد گویا باشد، کار
تمام است: می‌توانیم $a = b = \sqrt{2}$ بگذاریم. وگرنه این عدد گنگ
است. در این صورت داریم
\[
(\sqrt{2}^{\sqrt{2}})^{\sqrt{2}} = \sqrt{2}^{\sqrt{2} \cdot
  \sqrt{2}} = \sqrt{2}^2 = 2,
\]
که گویاست. پس در این حالت $a$ را~$\sqrt{2}^{\sqrt{2}}$ و~$b$ را
$\sqrt 2$ می‌گیریم.
\end{proof}

آیا این یک اثبات معتبر است؟ بیشتر ریاضی‌دانان آن را معتبر می‌دانند.
اما باز چیزی اندکی ناخوشایند در کار است: وجود یک جفت عدد حقیقی با
خاصیتی معین را ثابت کرده‌ایم، بی‌آنکه بتوانیم بگوییم \emph{کدام} جفت
عدد است. همان نتیجه را می‌توان به‌گونه‌ای ثابت کرد که جفت~$a$ و~$b$
در خود اثبات داده شود: بگذارید $a = \sqrt{3}$ و~$b = \log_3 4$.
عدد~$\sqrt3$ گنگ است. همچنین~$\log_3 4$ گنگ است؛ زیرا اگر
$\log_3 4=m/n$ برای اعداد طبیعی مثبت~$m,n$ می‌بود، آنگاه
$3^m=4^n=2^{2n}$، برخلاف یکتایی تجزیه به عوامل اول. بنابراین هر دو
$a$ و~$b$ گنگ‌اند و
\[
a^b = \sqrt{3}^{\log_3 4} = 3^{1/2 \cdot \log_3 4} = (3^{\log_3
  4})^{1/2} = 4^{1/2}= 2,
\]
زیرا~$3^{\log_3 x} = x$.

منطق شهودگرایانه برای صورت‌بندی گونه‌ای از استدلال طراحی شده است که
حرکت‌هایی مانند حرکت اثبات نخست را مجاز نمی‌داند. اثبات وجود~$x$ای که
$!A(x)$ را برآورده سازد یعنی باید~$x$ معینی عرضه کنید و، مانند اثبات
دوم، ثابت کنید که~$!A$ را برآورده می‌سازد. اثبات اینکه~$!A$ یا~$!B$
برقرار است مستلزم آن است که بتوانید یکی از آن دو را اثبات کنید.

به بیان صوری، منطق شهودگرایانه از محدودکردن دستگاه
!!a{derivation} منطق کلاسیک به شیوه‌ای معین حاصل می‌شود. از دیدگاه
ریاضی، این‌ها صرفاً دستگاه‌های استنتاجی صوری‌اند، اما چنان‌که گفته شد،
هدفشان صورت‌بندی نوعی استدلال ریاضی است. می‌توان آن را نوع استدلالی
دانست که بر مبنای دیدگاهی فلسفی دربارهٔ ریاضیات (مانند شهودگرایی
براور) موجه است؛ می‌توان آن را نوعی استدلال ریاضی «ملموس‌تر» و
رضایت‌بخش‌تر (همسو با سازنده‌گرایی بیشاپ) دانست؛ و می‌توان دربارهٔ اینکه
آیا توصیف صوری انگیزهٔ غیرصوری را به‌درستی صورت‌بندی می‌کند یا نه بحث
کرد. اما هر موضع فلسفی‌ای که داشته باشیم، می‌توانیم منطق شهودگرایانه را
به‌عنوان منطقی با ارائهٔ صوری مطالعه کنیم؛ و به هر دلیل، بسیاری از
منطق‌دانان ریاضی مطالعهٔ آن را جالب می‌یابند.

\end{document}
